\documentclass[11pt]{article}
\usepackage{docstyle}
\renewcommand\sheetname{Classwork \textperiodcentered{} Green Machines + Metals}
\begin{document}
\sheethead{Classwork}{Green Machines + Metallic materials \textperiodcentered{} Year 9 Science}{Each question: A / M / E}
\textbf{Print this sheet at 100\,\%} (the graph grid is true size).

\section{Part A \textperiodcentered{} Flowers and seeds}
\Q Complete the sentences.\lvl{A}
\par Pollination is the transfer of pollen from the \blank[24mm]{anther} to the \blank[24mm]{stigma}.
In \blank[22mm]{cross}-pollination the pollen comes from a different plant of the same species.
Fertilisation happens when the pollen \blank[24mm]{nucleus} fuses with the ovule nucleus, forming a \blank[24mm]{zygote}.\par\vspace{2mm}

\Q A dandelion seed has a parachute of fine hairs. Explain how this helps the dandelion plant.\lvl{M}
\alines{3}{The hairs catch the wind, so the seed is carried far away from the parent plant. The new plant then does not compete with the parent for light, water, minerals and space.}

\Q Give one advantage and one disadvantage of a plant reproducing asexually. Explain each.\lvl{E}
\alines{4}{Advantage: only one parent is needed, so no pollinator is needed, and the offspring are identical, so they keep the parent's good features. Disadvantage: there is no variation, so if the environment changes or a disease arrives, all the offspring could die.}

\section{Part B \textperiodcentered{} Photosynthesis and respiration}
\Q Complete the word equation for respiration.\lvl{A}
\par \blank[30mm]{glucose} + \blank[30mm]{oxygen} $\longrightarrow$ \blank[34mm]{carbon dioxide} + \blank[24mm]{water} (+ energy)\par\vspace{2mm}

\Q Before a photosynthesis experiment, a student keeps a plant in a dark cupboard for 2 days. A leaf tested with iodine then stays brown. Explain why she does this first.\lvl{M}
\alines{3}{In the dark the plant cannot photosynthesise and uses up the starch in its leaves. Any starch found at the end of the experiment must then have been made during the experiment.}

\Q Palisade cells are near the top of the leaf and are packed with chloroplasts. Explain how this helps the leaf.\lvl{E}
\alines{4}{The top of the leaf gets the most light. Chloroplasts contain chlorophyll, which absorbs light energy for photosynthesis. More chloroplasts where the light is strongest means more glucose can be made.}

\section{Part C \textperiodcentered{} Transport and osmosis}
\Q The table shows the height of the liquid in a Visking tubing model.\lvl{M}
\par\vspace{1mm}
\begin{tabular}{@{}l *{8}{c}@{}}
\toprule
Time (min) & 0 & 5 & 10 & 15 & 20 & 25 & 30 & 35\\
Height (mm) & 10 & 22 & 34 & 46 & 58 & 70 & 72 & 73\\
\bottomrule
\end{tabular}\par\vspace{2mm}
\sub{a} Plot the results on the grid. Label both axes with units, choose an even scale, and draw a line.
\par\vspace{1mm}
\hbox to\linewidth{\hss\GraphGrid{14}{8}{%
  \foreach \t/\h in {0/10,5/22,10/34,15/46,20/58,25/70,30/72,35/73}
    \draw (\t*0.4-0.1,\h*0.1-0.1) -- (\t*0.4+0.1,\h*0.1+0.1) (\t*0.4-0.1,\h*0.1+0.1) -- (\t*0.4+0.1,\h*0.1-0.1);
  \draw (0,1.0) -- (10,7.0) -- (12,7.2) -- (14,7.3);
  \foreach \t in {0,5,...,35} \node[below] at (\t*0.4,0) {\small\t};
  \foreach \h in {0,20,...,80} \node[left] at (0,\h*0.1) {\small\h};
  \node[below=4mm] at (7,0) {\small time (min)};
  \node[rotate=90, above=6mm] at (0,4) {\small height (mm)};}\hss}
\ansnote{5mm}{Scale: 1 cm = 2.5 min (x), 1 cm = 10 mm (y). Straight line from 0 to 25 min, then nearly flat.}
\sub{b} Describe the pattern in the results.\lvl{M}
\alines{2}{The height rises steadily, by 12 mm every 5 minutes, up to 25 minutes. After 25 minutes it hardly rises at all.}
\sub{c} Suggest why the liquid stops rising.\lvl{E}
\alines{2}{Water has moved in, so the sugar solution is more dilute. The difference in water concentration is smaller, so less water moves in by osmosis.}

\section{Part D \textperiodcentered{} Properties and density}
\Q Name the property of metals that each sentence shows.\lvl{A}
\par\sub{a} Copper can be drawn out into thin wire. \hfill\blank[34mm]{ductile}
\par\sub{b} Gold can be beaten into very thin sheets. \hfill\blank[34mm]{malleable}
\par\sub{c} A brass bell rings when it is struck. \hfill\blank[34mm]{sonorous}\par\vspace{1mm}

\Q A metal cube has sides of 3.0\,cm and a mass of 72.9\,g. Calculate its density and use the booklet table (iron 7.9, copper 8.9, silver 10.5, aluminium 2.7\,g/cm$^{3}$) to identify the metal. Show your working.\lvl{M}
\flines{3}{volume $= 3.0 \times 3.0 \times 3.0 = 27.0\ \mathrm{cm^{3}}$\\
density $= \dfrac{72.9\ \mathrm{g}}{27.0\ \mathrm{cm^{3}}} = 2.7$ g/cm$^{3}$\\
The metal is aluminium.}

\section{Part E \textperiodcentered{} Alloys and heat treatment}
\Q Complete: an alloy is a mixture of two or more elements, at least one of which is a \blank[30mm]{metal}.\lvl{A}

\Q Steel is iron with a small amount of carbon. Explain why steel is harder than pure iron. Use the idea of layers of atoms.\lvl{E}
\alines{4}{In pure iron the atoms are all the same size and sit in regular layers, which slide over each other easily. Carbon atoms are a different size, so they disrupt the layers. The layers cannot slide easily, so steel is harder.}

\Q Glasses frames made of nitinol can be bent and still go back to their shape. What makes nitinol a smart alloy?\lvl{M}
\alines{2}{It is a shape memory alloy: after it is bent out of shape, it returns to its original shape when it is heated (or a current passes through it).}

\section{Part F \textperiodcentered{} Corrosion and reactions}
\Q Write the word equations.\lvl{A}
\par\sub{a} calcium + hydrochloric acid $\longrightarrow$ \blank[78mm]{calcium chloride + hydrogen}
\par\sub{b} sodium + oxygen $\longrightarrow$ \blank[78mm]{sodium oxide}\par\vspace{1mm}

\Q Plan an experiment to put magnesium, iron and copper in order of reactivity, using dilute hydrochloric acid. Include your method, how you make it a fair test, what you would observe, your prediction and one safety step.\lvl{E}
\alines{8}{Method: put the same volume (5 mL) of the same dilute hydrochloric acid into three test tubes. Add a piece of each metal of the same size and shape (cleaned with sandpaper). Fair test: same volume and concentration of acid, same size of metal, same temperature. Observe: how fast bubbles of hydrogen are given off. Prediction: magnesium bubbles fastest, iron slowly, copper not at all, so the order is magnesium, iron, copper. Safety: wear safety glasses, because acid can damage eyes.}
\end{document}
